\documentclass[11pt]{article}

\usepackage[margin=1in]{geometry}
\usepackage{amsmath,amssymb,booktabs,tabularx}

\title{First-moment guard for \(Z_t\)}
\author{}
\date{}

\begin{document}
\maketitle

\noindent\textbf{Status.} Draft\\
\textbf{Plan item.} First-moment guard for \(Z_t\)\\
\textbf{Depends on.} \texttt{01\_long\_variance\_normalization.tex}\\
\textbf{Feeds.} P13, affine domain constraints, calibration admissibility, pricing admissibility

\section{Purpose}

Derive the admissibility condition that keeps the unit-mean long-variance
normalization compatible with a positive mean-reversion level.
The guard prevents uncompensated event feedback from exhausting the entire
stationary first-moment budget of \(Z_t\).

\section{Setup}

Use the physical long-variance dynamics
\[
    dZ_t
    =
    \kappa_z(\bar z-Z_t)\,dt
    +\sigma_z\sqrt{Z_t}\,dW_t^Z
    +\chi_+\,dN_t^+
    +\chi_-\,dN_t^-,
\]
with stationary event-feedback drift
\[
    b_Z
    =
    \chi_+\lambda_{\mathrm{stat}}^+
    +\chi_-\lambda_{\mathrm{stat}}^-.
\]

From the unit-mean normalization,
\[
    \bar z=1-\frac{b_Z}{\kappa_z}.
\]

\section{Notation}

\begin{tabularx}{\textwidth}{lXX}
\toprule
Symbol & Meaning & Units or domain\\
\midrule
\(\kappa_z\) & Mean-reversion rate & Positive inverse time\\
\(\bar z\) & Mean-reversion level & Positive\\
\(\chi_+\) & Buy-event long-variance jump loading & Non-negative\\
\(\chi_-\) & Sell-event long-variance jump loading & Non-negative\\
\(\lambda_{\mathrm{stat}}^+\) & Stationary buy event rate & Non-negative inverse time\\
\(\lambda_{\mathrm{stat}}^-\) & Stationary sell event rate & Non-negative inverse time\\
\(b_Z\) & Stationary event-feedback drift into \(Z_t\) & Non-negative inverse time\\
\bottomrule
\end{tabularx}

\section{Assumptions used}

\begin{tabularx}{\textwidth}{XXX}
\toprule
Assumption & Where it is used & Why it is needed\\
\midrule
The unit-mean normalization fixes \(\bar z=1-b_Z/\kappa_z\).
& Rewriting positivity of \(\bar z\).
& The guard is the positivity condition for the derived level.\\
The long-variance mean-reversion level is strictly positive.
& Applying \(\bar z>0\).
& The CIR-type drift target must remain an admissible positive level.\\
\(\kappa_z>0\).
& Multiplying the inequality by \(\kappa_z\).
& The inequality direction is preserved and the process mean-reverts.\\
Event loadings and stationary event rates are non-negative.
& Interpreting \(b_Z\) as consumed first-moment budget.
& Positive event feedback can only reduce the retained level under unit normalization.\\
\bottomrule
\end{tabularx}

\section{Result}

\noindent\textbf{Proposition.} Under the unit-mean normalization, the
long-variance first moment is admissible only if
\[
    \kappa_z
    >
    \chi_+\lambda_{\mathrm{stat}}^+
    +\chi_-\lambda_{\mathrm{stat}}^-.
\]

Equivalently,
\[
    b_Z<\kappa_z.
\]

\section{Derivation}

The unit-mean normalization gives
\[
    \bar z
    =
    1-\frac{b_Z}{\kappa_z}.
\]

The long-variance drift target must satisfy \(\bar z>0\). Substituting the
normalized level,
\[
    1-\frac{b_Z}{\kappa_z}>0.
\]

Because \(\kappa_z>0\), multiply by \(\kappa_z\) without changing the inequality
direction:
\[
    \kappa_z-b_Z>0.
\]

Therefore,
\[
    \kappa_z>b_Z.
\]

Replacing \(b_Z\) by the event-feedback drift gives
\[
    \kappa_z
    >
    \chi_+\lambda_{\mathrm{stat}}^+
    +\chi_-\lambda_{\mathrm{stat}}^-.
\]

The boundary case \(b_Z=\kappa_z\) gives \(\bar z=0\), which violates the
strict positive-level contract.
The supercritical case \(b_Z>\kappa_z\) gives \(\bar z<0\), so the unit-mean
target would require a negative mean-reversion level.

\section{Domain and validity}

The guard is valid under the same first-moment conditions as the normalization:
\[
    \kappa_z>0,\qquad
    \chi_+\ge 0,\qquad
    \chi_-\ge 0,\qquad
    \lambda_{\mathrm{stat}}^+\ge 0,\qquad
    \lambda_{\mathrm{stat}}^-\ge 0,
\]
with finite stationary side event rates.
If the Hawkes system is not stationary, or if the stationary event rates are
not finite, then \(b_Z\) is not a valid finite first-moment input and the guard
cannot be evaluated as a finite inequality.

\section{Consequences}

\begin{itemize}
    \item The event-feedback drift consumes part, but not all, of the long-variance first-moment budget.
    \item The retained mean-reversion level \(\bar z\) remains positive.
    \item Calibration must project or reject event loadings whose stationary event-feedback drift reaches or exceeds \(\kappa_z\).
    \item Pricing and affine-transform inputs must use parameter sets that satisfy the same first-moment admissibility condition.
\end{itemize}

\section{Checks}

\begin{tabularx}{\textwidth}{XXX}
\toprule
Check & Expected result & Evidence location\\
\midrule
Exact guard substitution
& \(b_Z<\kappa_z\) if and only if the derived \(\bar z\) is positive.
& \texttt{check\_long\_variance\_first\_moment.py}\\
Kernel guard contract
& Parameter sets with \(b_Z\ge\kappa_z\) fail stationary-level validation.
& Unit test \texttt{stationary\_mean\_reversion\_level\_rejects\_exhausted\_event\_feedback\_budget}\\
\bottomrule
\end{tabularx}

\end{document}
